\subsection{Applications.}

* A) Unitary representations of groups :

Let E be a complex hilbertian space, G a group and \(\pi\) a unitary representation of G on E, i.e.~a homomorphism from G into the group of automorphisms of E. Let \(E^{G}\) be the hilbertian subspace of E consisting of all vectors invariant under \(\pi(\mathbf{G})\) . For every \(x \in E\) , let \(K_{x}\) be the closed convex envelope of the orbit of x. Fix a point x in E.

We shall show that there exists a unique point in \(K_{x}\) which is invariant under \(\pi(\mathbf{G})\) , namely the projection of x on \(E^{G}\) . By IV, p.~44, corollary (applied to the underlying real vector space of E), there exists a point in \(K_{x}\) which is invariant under \(\pi(\mathbf{G})\) ; let a be such a point; then \(a \in E^{G}\) . Let P be the set of all \(y \in E\) such that y - x is orthogonal to \(E^{G}\) ; we see immediately that P is closed, convex and invariant under \(\pi(\mathbf{G})\) ; therefore \(x \in P\) , hence \(K_{x} \subset P\) and finally \(a \in P\) . In other words, a - x is orthogonal to \(E^{G}\) ; consequently a is the projection of x onto \(E^{G}\) .

* B) Trace of an operator in a hilbertian space :

Suppose that the representation \(\pi\) is irreducible, that is, that there exists no hilbertian subspace of E, distinct from \(\{0\}\) and from E, which is invariant under \(\pi(\mathbf{G})\) . Let \(\mathbf{F} = \mathcal{L}^{2}(\mathbf{E})\) be the hilbertian space of all Hilbert-Schmidt endomorphisms of E, with the scalar product \(\langle u|v\rangle = \operatorname{Tr}(u^{*}v)\) . We define a unitary representation \(\lambda\) from G into F by the formula

\[
\lambda (g). u = \pi (g) u \pi (g) ^ {- 1} \quad (u \in \mathrm{F}, g \in \mathrm{G}).\tag{3}
\]

The space \(F^{G}\) of all elements of E invariant under \(\lambda(\mathbf{G})\) consists of the Hilbert-Schmidt endomorphisms u of E which commute with \(\pi(g)\) for all \(g \in G\) . By Schur's lemma, such a u is a homothety. Hence we must consider two cases:

\begin{enumerate}
\def\labelenumi{\arabic{enumi})}
\item
  if E is infinite dimensional, then \(F^{G} = \{0\}\) ;
\item
  if \(\mathrm{E}\) is finite dimensional, then \(\mathrm{F} = \mathcal{L}(\mathrm{E})\) and \(\mathrm{F}^{\mathrm{G}} = \mathbf{C}.1_{\mathrm{E}}\) .
\end{enumerate}

Applying the result of \(A\) ) to the unitary representation \(\lambda\) , we obtain the following theorem :

Let \(u \in \mathcal{L}^2(\mathrm{E})\) , and let \(\mathbf{A}_u\) be the closed convex envelope in \(\mathcal{L}^2(\mathrm{E})\) of the set of endomorphisms \(\pi(g) u\pi(g)^{-1}\) of \(\mathrm{E}\) , where \(g\) runs through \(\mathbf{G}\) . If \(\mathrm{E}\) is infinite dimensional, we have \(0 \in \mathbf{A}_u\) . If \(\mathrm{E}\) is finite dimensional with dimension \(d\) , there exists a unique homothety

in \(\mathbf{A}_u\) , namely the projection \(\frac{1}{d}\mathrm{Tr}(u).1_{\mathrm{E}}\) of \(u\) onto the subspace \(\mathbf{C}.1_{\mathrm{E}}\) of \(\mathcal{L}^2 (\mathrm{E}).\)

\begin{enumerate}
\def\labelenumi{\Alph{enumi})}
\setcounter{enumi}{2}
\tightlist
\item
  Haar measure of a compact group :
\end{enumerate}

Let G be a compact group and let \(E = \mathcal{C}(\mathbf{G}, \mathbf{R})\) be the Banach space of all real valued continuous functions on G, endowed with the norm

\[
\| f \| = \sup _ {x \in \mathbf {G}} | f (x) |.\tag{4}
\]

For all \(x \in G\) , we define the automorphisms \(\gamma_{x}\) and \(\delta_{x}\) of E by the formulas

\[
\gamma_ {x} f (y) = f \left(x ^ {- 1} y\right), \quad \delta_ {x} f (y) = f (y x)\tag{5}
\]

(for \(y \in \mathbf{G}\) , \(f \in \mathrm{E}\) ).

Let \(f \in E\) , let \(\Gamma_f\) (resp. \(\Delta_f\) ) denote the closed convex envelope in E of the set of all functions \(\gamma_x f\) (resp. \(\delta_x f\) ) as \(x\) ranges over G. We shall prove that there exists a unique constant function \(\mu(f)\) belonging to \(\Gamma_f\) , a unique constant function \(\mu'(f)\) belonging to \(\Delta_f\) , and that these constants are equal.

It is clear that a continuous function on G is invariant under the automorphisms \(\gamma_{x}(\text{resp. } \delta_{x})\) of E if and only if it is constant. Now the set of all functions \(\gamma_{x}f(\text{resp. } \delta_{x}f)\) for x in G, is compact in E, since the mapping \(x \mapsto \gamma_{x}f(\text{resp. } x \mapsto \delta_{x}f)\) from G into E is continuous (GT, X, §3, No.~4, th. 3). It follows (II, p.~25, prop. 3) that \(\Gamma_{f}(\text{resp. } \Delta_{f})\) is a compact set in E for the topology defined by the norm, hence for \(\sigma(\text{E}, \text{E}')\) . By the Ryll-Nardzewski th. (IV, p.~43, th. 2), there exist constant functions in \(\Gamma_{f}\) and \(\Delta_{f}\) . It remains to prove that if \(c_{1} \in \Gamma_{f}\) and \(c_{2} \in \Delta_{f}\) are constants, then \(c_{1} = c_{2}\) .

Let \(\varepsilon > 0\) . By the hypothesis there exist points \(x_{1}, \ldots, x_{n}, y_{1}, \ldots, y_{n}\) in \(\mathbf{G}\) and positive real numbers \(\lambda_{1}, \ldots, \lambda_{n}, \mu_{1}, \ldots, \mu_{m}\) such that

\[
\lambda_ {1} + \dots + \lambda_ {n} = \mu_ {1} + \dots + \mu_ {m} = 1.\tag{6}
\]

\[
\sup _ {x \in \mathbf {G}} \left| \sum_ {i = 1} ^ {n} \lambda_ {i} f (x _ {i} x) - c _ {1} \right| \leqslant \varepsilon ,\tag{7}
\]

\[
\sup _ {x \in \mathbf {G}} \left| \sum_ {j = 1} ^ {m} \mu_ {j} f (x y _ {j}) - c _ {2} \right| \leqslant \varepsilon .\tag{8}
\]

Put \(r = \sum_{i,j} \lambda_i \mu_j f(x_i y_j)\) . Then \(r - c_1 = \sum_{j=1}^{m} \mu_j a_j\) with \(a_j = \sum_{i=1}^{n} \lambda_i f(x_i y_j) - c_1\) ; by (7), we have \(|a_j| \leqslant \varepsilon\) for \(1 \leqslant j \leqslant m\) , hence \(|r - c_1| \leqslant \varepsilon\) . Similarly, we prove the inequality \(|r - c_2| \leqslant \varepsilon\) , hence \(|c_1 - c_2| \leqslant 2\varepsilon\) . Since \(\varepsilon\) is arbitrary, we get \(c_1 = c_2\) , as asserted.

By the definition of \(\mu(f)\) , for every \(\varepsilon > 0\) we can find positive numbers \(\lambda_1, \ldots, \lambda_n\) with sum 1 and elements \(x_1, \ldots, x_n\) in \(\mathbf{G}\) such that \(\left|\sum_{i=1}^{n} \lambda_i f(x_i x) - \mu(f)\right| \leqslant \varepsilon\) for all \(x \in \mathbf{G}\) .

It is immediate that for \(f, g\) in \(\mathbf{E}\) and for every scalar \(\lambda\) , we have \(\Gamma_{f + g} \subset \Gamma_f + \Gamma_g\) and \(\Gamma_{\gamma f} = \lambda \Gamma_f\) , hence we have the relations \(\mu(f + g) = \mu(f) + \mu(g)\) and \(\mu(\lambda f) = \lambda \mu(f)\) . Therefore, the mapping \(\mu : f \mapsto \mu(f)\) from \(\mathbf{E}\) into \(\mathbf{R}\) is a mean on the compact space \(\mathbf{G}\) (IV, p.~40); * in other words \(\mu\) is a positive measure on \(\mathbf{G}\) such that \(\mu(\mathbf{G}) = 1\) . *

It is immediate that \(\mu\) is invariant under the left translations of G, and the equality \(\mu(f) = \mu'(f)\) implies that \(\mu\) is also invariant under right translations. * In other words, \(\mu\) is a left and a right measure on G (INT, VII, § 1, No.~2, def. 2). *

\paragraph{* D) Existence of invariant measures :}

Let X be a Hausdorff topological space, \(\mu\) a positive bounded measure on X, and G a group of homeomorphisms of X. Suppose that for all \(g \in G\) , the measure \(g \cdot \mu\) , the image of \(\mu\) under the mapping \(g : X \to X\) is of base \(\mu\) . Let \(u_g\) be a positive \(\mu\) -integrable function on X such that \(g \cdot \mu = u_g \cdot \mu\) . Suppose also that there exist two positive \(\mu\) integrable functions \(\phi\) and \(\psi\) on X, which are not \(\mu\) -null and are such that \(\phi \leqslant u_g \leqslant \psi\) \(\mu\) -almost everywhere for all \(g \in G\) . We shall prove that there exists a positive bounded measure \(v \neq 0\) on X, with base \(\mu\) , and invariant under G.

Let P be the subset of the Banach space \(E = L^{1}(X, \mu)\) consisting of classes of functions f such that \(\phi \leqslant f \leqslant \psi\) \(\mu\) -almost everywhere. Then P is compact for the weakened topology \(\sigma(E, E')\) . The mapping \(h \mapsto h \cdot \mu\) from P into the Banach space \(F = \mathcal{M}^{b}(X)\) of bounded real measures on X, is a bijection from P onto a subset \(P_{1}\) of E which is convex and compact for the topology \(\sigma(F, F')\) . By hypothesis, \(g \cdot \mu \in P_{1}\) for all \(g \in G\) . Let K be the closed convex envelope of the set of all measures \(g \cdot \mu\) . For all \(g \in G\) , the mapping \(v \mapsto g \cdot v\) is an isometric affine transformation of K. By the Ryll-Nardzewski th. (IV, p.~43, th. 2), there exists a measure \(v \in K\) which is invariant under G. We have \(\phi \cdot \mu \leqslant v\) , hence \(v \neq 0\) .

\closechapterappendix
